\chapter{生成式AI對教育體系的衝擊}
\label{chap:生成式AI對教育體系的衝擊}

\chapterauthor{陳慶霖、鍾詠傑、林吟貞、林駿甫}

\section{摘要}\label{sec:摘要}

本報告旨在從不同角度了解AI對教育體系之影響，(一)各國指引：此部分整理現有各國政府教育部門釋出的AI融入教學指引，並進一步進行比較。(二)教育現場的調查：第二部分對台灣教育現場的師生進行初步調查，以問卷蒐集120位國中九年級學生及16位中小學教師對AI融入教學態度與想法，進一步分析與討論。(三)文獻回顧：透過回顧AI融入教學的實徵研究，來了解對學習所帶來的影響有哪些。(四)綜合觀點：在這些結果中我們如何看待AI是否應融入教學，並進一步提出未來建議。

\section{議題簡介：時事案例}\label{sec:時事案例}

二〇二六年四月，美國普渡大學（Purdue University）資訊工程學系爆發了一起震驚全美學術界的事件。該系核心必修課程「CS240 系統程式設計（Systems Programming）」的授課教師 Jeffrey Turkstra 副教授，透過其參與開發並於 2022 年發表於電腦科學教育領域頂尖國際會議 SIGCSE 的程式碼追蹤系統「EnCourse」\cite{rodriguez2022tracking}，向班上超過兩百名學生發出學術誠信違規（Academic Integrity Violation）的指控電子郵件\cite{arun2026ai,lafayette2026cheating}。

這封電子郵件的時機與內容引發了極大的爭議。Turkstra 選擇在該學期的「退選截止日（Withdrawal Deadline）」當天寄出此信，要求被指控的學生填寫一份 Google 表單，自行坦白在哪幾份作業中使用了生成式 AI 工具，並警告若未於限期內回覆，可能面臨課程成績不及格（F）及學務處（Office of the Dean of Students）的進一步調查與懲處\cite{arun2026ai,lafayette2026cheating}。更引發學生恐慌的是，信中並未告知每位學生具體是哪一份作業觸發了系統的警報，學生必須在缺乏具體資訊的情況下自行認罪，形同將舉證責任完全轉嫁給學生\cite{arun2026ai,bailey2026cheating}。

郵件在學生社群中引發了大規模恐慌。根據《Lafayette Journal \& Courier》（USA TODAY Network）的報導，超過半數的學生選擇在截止日前緊急退選課程\cite{lafayette2026cheating}。消息迅速在網路社群擴散，累積了數千則討論，引起了校園報紙《The Purdue Exponent》及主流新聞媒體的關注\cite{arun2026ai,lafayette2026cheating}。在強烈的輿論壓力與學務處的行政介入下，Turkstra 於四月二十日的課堂上宣布了全面的政策撤回：過往所有作業的 AI 作弊指控一律撤銷，認罪表單收集的資料全數作廢，已退選的學生可無條件重新加選\cite{arun2026ai}。

Turkstra 在課堂上展示了課程數據，指出使用 AI 完成作業的學生在實體閉卷考試中的成績，比親自完成作業的學生平均低了百分之十至十五，佐證其在基礎課程中限制 AI 使用的教學立場\cite{arun2026ai}。然而，全面撤銷指控也帶來了新的道德困境。Turkstra 坦言：「我們有許多投入大量心血、沒有作弊的學生，現在卻只能眼睜睜看著部分作弊的同學安然脫身且不用承擔任何後果。對此，我無法給出一個滿意的答案。」\cite{arun2026ai}

這起事件精準地折射出生成式 AI 對教育體系帶來的多重衝擊：在技術層面，自動化偵測工具的準確性與誤判問題亟待解決；在教學層面，基礎課程中應否禁止使用 AI 工具、以及如何重新設計評量方式，成為教育工作者必須面對的核心課題；在制度層面，程序正義與學術誠信之間的平衡，在 AI 時代面臨前所未有的挑戰。這些問題，也正是本報告希望從不同角度深入探討的焦點。

\section{現況分析：政府的指引}\label{sec:政府指引}

面對生成式 AI 帶來的教育衝擊，各國政府與國際組織紛紛制定政策指引，試圖在擁抱科技創新與維護學術品質之間取得平衡。聯合國教科文組織（UNESCO）於二〇二四年正式發布了《學生人工智慧素養架構》（AI Competency Framework for Students），為全球教育體系提供了首份標準化的 AI 能力參考藍圖\cite{unesco2024competency}。根據 UNESCO 的調查，截至二〇二二年，全球僅有十一個國家正式制定並核定 K-12 AI 課程綱要，另有四個國家正在研發中\cite{unesco2022k12}，凸顯了教育體系採用新技術的速度遠落後於科技發展，已成為全球性的教育危機。UNESCO 的素養架構並未將 AI 教育窄化為單純的程式設計或軟體操作教學，而是建構了一套由「四個核心維度」與「三個漸進階層」交織而成的十二項素養矩陣\cite{unesco2024competency}。四個維度涵蓋了「以人為本的心智模式」、「人工智慧倫理」、「AI 技術與應用」以及「AI 系統設計」；三個漸進階層則從「理解（Understand）」出發，經由「應用（Apply）」，最終達到「創造（Create）」的精熟水準。架構的核心精神在於強調人類主體性（Human Agency）與當責性（Accountability）：即使 AI 系統提供了看似權威的建議，人類依然擁有最終的決策權，並必須為 AI 產出的內容承擔完全責任\cite{unesco2024competency}。

在國內方面，臺灣教育部因應生成式 AI 的快速發展，將《中小學使用生成式人工智慧注意事項》更新至 2.1 版，並依據不同學習階段的認知發展特徵，編製了兩套學習應用手冊：針對國小學生的《我和 AI 一起學》以及針對國高中學生的《駕馭 AI，洞察未來：數位公民的必修課》\cite{govtw2024aiguide}。國小版手冊以「5W1H」架構引導學童建立對 AI 的基礎認知與自我保護意識，強調不洩露個人隱私資料、辨識 AI 幻覺（Hallucination）等核心觀念；國高中版則將教育目標提升至培養「AI 共學力」，強調操作技能精進、批判性內容查核與倫理反思\cite{govtw2024aiguide}。在高等教育層面，各大專校院也紛紛發布校內指引。以國立臺灣師範大學為例，其教學發展中心公布的《生成式 AI 之學習應用及參考指引》明確列出五大核心倫理面向（公平性、隱私與資料保護、透明與可解釋性、責任歸屬以及人機界線與原創性），並要求師生在使用 AI 時必須進行自我揭露，清楚標註 AI 參與的範疇，同時鼓勵各教學單位依學門特性自訂更細緻的規範\cite{ntnu2025ai}。

跨國比較則揭示了各國在政策取向上的顯著差異。日本文部科學省於二〇二四年十二月發布了《初等中等教育階段生成式 AI 使用指引（Ver. 2.0）》（初等中等教育段階における生成AIの利活用に関するガイドライン），展現了最為審慎的立場\cite{mext2024genai}。該指引強調不應統一禁止或強制使用 AI，而是應評估其使用是否有助於課程綱要所揭示的素養與能力培育，但同時也明確警告，過度依賴 AI 可能導致學生對 AI 生成答案的「盲目接受」，進而對「思考力、判斷力與表達力」的發展造成不利影響\cite{mext2024genai}。

相較之下，英國教育部（DfE）的 \textit{Generative Artificial Intelligence (AI) in Education} 政策文件則採取了「雙軌制」策略：一方面大力鼓勵教師端使用 AI 以減輕行政負擔（如課程規劃、回饋生成），甚至資助 Oak National Academy 開發了 AI 教學助理工具「Aila」\cite{dfe2023genai}，另一方面則要求學校在使用任何 AI 工具前必須進行全面的風險評估，嚴格遵守資料保護法與《在學校中維護兒童安全》（\textit{Keeping Children Safe in Education}）的法定責任，並關注深度偽造（Deepfakes）等新興風險\cite{dfe2023genai}。

至於美國，教育部教育技術辦公室於二〇二三年五月發布的 \textit{Artificial Intelligence and the Future of Teaching and Learning} 報告則將視角拉高至社會公平層面\cite{usedu2023ai}。該報告明確提出 AI 應「增強人類智能而非取代人類」，強調必須保持「人類在迴路（Human-in-the-loop）」的核心原則，並聚焦管理「演算法偏見」的風險，因為 AI 系統的開發過程可能導致模式偵測的偏差與自動化決策的不公，進而邊緣化弱勢群體\cite{usedu2023ai}。

新加坡則展現了最為積極的擁抱姿態。根據經濟合作暨發展組織（OECD）二〇二四年《教與學國際調查》（TALIS）的數據，新加坡高達百分之七十五的中學教師已在教學中常態性使用 AI，遠超 OECD 百分之三十六的平均水準\cite{moe2025singapore}。新加坡教育部提出了「高信任、高接觸、高科技」的教育架構，將 AI 定位為實現「有效班級規模為一」（即高度個人化學習）的關鍵工具，同時堅持人類教師在社會情感教育中的不可替代性\cite{moe2025singapore}。在高教端，新加坡國立大學（NUS）制定了嚴謹的生成式 AI 使用規範，確立了「使用者須對 AI 生成的所有內容負完全責任，而非 AI 工具本身」的核心原則，並要求學生在提交作業時必須揭露使用的 AI 工具名稱、使用方式及具體貢獻範圍\cite{nus2024integrity}。

綜觀各國政策，可以發現若干高度一致的共識：第一，AI 不應取代人類的判斷與決策，教育者與學習者必須保持「人類在迴路」的主體地位；第二，所有允許 AI 參與的場景，皆以「強制揭露」與「人類負完全責任」為前提；第三，各國普遍認知到，教育體系不應僅停留在「防堵 AI」的消極立場，而應積極建構系統性的 AI 素養教育。然而，在開放程度與監管強度上，各國仍有顯著分歧，這也反映了不同文化底蘊、教育體制與國家科技發展戰略的深層差異。這些官方指引為教育現場提供了原則性架構，但實際的執行成效與師生的真實態度，仍需要透過實證調查來進一步驗證。

\section{實際調查：問卷設計與發放}\label{sec:問卷設計}

本報告旨在探討教學現場師生對生成式人工智慧（AI）融入教育的看法，因此採用問卷調查法進行資料蒐集。根據生成式AI於教育領域的應用現況及相關文獻資料，自行設計教師版與學生版兩份問卷，以了解師生對AI工具的使用情形、態度及影響認知。

本研究以\textbf{120位國中九年級學生}及\textbf{16位中小學教師}為調查對象，屬初步探索性調查。學生問卷主要針對學生在學習過程中使用生成式AI工具的情形進行調查；教師問卷則聚焦於教師在教學現場使用AI工具的經驗，以及對AI融入教育的看法。問卷採匿名方式填寫，以提高受試者作答意願及資料真實性。

學生問卷內容包含五個面向：（一）AI工具使用頻率與情境、（二）AI使用方式、（三）AI依賴程度、（四）學習行為是否改變、（五）AI倫理與使用態度，共24題。問卷採一至五分李克特尺度，分數越高代表越同意該題敘述；其中 Q1 至 Q4 主要用於了解基本使用經驗與頻率，後續量化構面分析則聚焦於 Q5 至 Q24。透過這些題項了解學生使用AI工具的實際情況及其對學習產生的影響。

教師問卷則包含背景資料調查、AI使用經驗、學生AI使用觀察、教師對AI融入教育的態度，以及開放式意見回饋等部分。內容涵蓋教師使用AI進行備課、教材設計、行政工作等情形，並探討教師對AI對學生學習、教學方式及未來教育發展之看法。

問卷回收後，研究團隊將資料進行整理與統計分析，藉以了解師生對AI融入教學的認知、使用現況及態度差異，作為未來教育現場推動AI應用之參考依據。不過，由於樣本來源與數量有限，且未採隨機抽樣，本調查結果不宜推論為全臺中小學師生的整體態度，而應作為觀察教育現場初步傾向的參考。

\section{分析結果}\label{sec:分析結果}

\subsection{學生的態度}\label{sec:學生態度}

\subsubsection{1. 構面：使用方式}

\begin{longtable}{l c c}
\toprule
題目 & 平均數 & 標準差 \\
\midrule
\endfirsthead
\toprule
題目 & 平均數 & 標準差 \\
\midrule
\endhead
\midrule
\multicolumn{3}{r}{續下頁} \\
\endfoot
\bottomrule
\endlastfoot
5. 我會先自己思考，再詢問AI & 4.00 & 1.000 \\
6. 我會直接使用 AI 提供的答案 & 2.68 & 1.033 \\
7. 我會修改 AI 產生的內容 & 3.66 & 1.203 \\
8. 我會確認 AI 回答是否正確 & 3.75 & 1.147 \\
9. 我會請 AI 解釋不懂的概念 & 4.09 & 0.974 \\
10. 我會透過反覆提問讓 AI 提供更好的答案 & 3.72 & 1.221 \\
\end{longtable}

在使用方式構面中，可以從結果分析看出「先自己思考再詢問 AI」（Q5，M = 4.00）以及「請 AI 解釋不懂的概念」（Q9，M = 4.09）兩題平均數較高，顯示\textbf{學生較常將 AI 用於概念理解與學習輔助}。

\subsubsection{2. 構面：依賴程度}

\begin{longtable}{l c c}
\toprule
題目 & 平均數 & 標準差 \\
\midrule
\endfirsthead
\toprule
題目 & 平均數 & 標準差 \\
\midrule
\endhead
\midrule
\multicolumn{3}{r}{續下頁} \\
\endfoot
\bottomrule
\endlastfoot
11. 沒有 AI 時，我會不知道如何開始作業 & 1.61 & 0.922 \\
12. 我習慣先詢問 AI，而不是自己查資料 & 2.22 & 1.206 \\
13. 我認為 AI 比自己搜尋資料更方便 & 3.17 & 1.174 \\
14. 使用 AI 後，我較少主動思考問題 & 2.13 & 1.094 \\
15. 我越來越依賴 AI 協助學習 & 2.18 & 1.110 \\
\end{longtable}

在依賴程度構面中，各題的平均數普遍偏低，尤其是 Q11 的分數接近極小值。反映出多數學生並未對 AI 產生過度依賴，仍自認具備獨立規劃並完成學習任務的核心能力。

\subsubsection{3. 構面：學習行為改變}

\begin{longtable}{l c c}
\toprule
題目 & 平均數 & 標準差 \\
\midrule
\endfirsthead
\toprule
題目 & 平均數 & 標準差 \\
\midrule
\endhead
\midrule
\multicolumn{3}{r}{續下頁} \\
\endfoot
\bottomrule
\endlastfoot
16. AI 改變了我的學習方式 & 2.65 & 1.222 \\
17. 使用 AI 後，我更有效率地完成學習任務 & 3.46 & 1.023 \\
18. 使用 AI 後，我花較少時間查資料 & 3.24 & 1.133 \\
19. 使用 AI 後，我更願意自主學習 & 2.92 & 1.005 \\
20. AI 讓我更容易理解困難內容 & 3.29 & 1.075 \\
\end{longtable}

在此構面中，Q17 與 Q20 的平均數落在 3.2 至 3.5 之間，屬於中立區間。這顯示學生對於「AI 是否顯著改變其學習效率或理解困難內容」的看法尚不一致，整體態度偏向中立，並未形成壓倒性的普遍認同。

\subsubsection{4. 構面：倫理與使用態度}

\begin{longtable}{l c c}
\toprule
題目 & 平均數 & 標準差 \\
\midrule
\endfirsthead
\toprule
題目 & 平均數 & 標準差 \\
\midrule
\endhead
\midrule
\multicolumn{3}{r}{續下頁} \\
\endfoot
\bottomrule
\endlastfoot
21. 我認為直接繳交 AI 生成內容是不恰當的 & 3.67 & 1.288 \\
22. 我會擔心 AI 提供錯誤資訊 & 4.10 & 0.982 \\
23. 我認為學生應學習如何正確使用 AI & 4.18 & 0.908 \\
24. 我認為 AI 對學習的幫助大於壞處 & 3.45 & 0.963 \\
\end{longtable}

在倫理與使用態度構面中，可以從題目「學生普遍認為應學習如何正確使用 AI」（Q23，M = 4.18），並對「AI 可能產生錯誤資訊抱持一定程度的擔憂」（Q22，M = 4.10），表示學生普遍認為\textbf{AI 仍應作為輔助工具，而不是直接取代自己的學習成果}。

\subsubsection{結論}

整體而言，可以從表單結果看出學生傾向肯定 AI 在概念理解、完成學習任務以及查詢資料上的輔助價值，但在「是否顯著改變學習效率或學習方式」上仍偏向中立，尚不能解讀為壓倒性的普遍認同。

此外，多數學生表示在使用 AI 前仍會先自行思考，並會確認或調整 AI 所提供的內容，顯示學生傾向將 AI 作為學習輔助工具，而非直接依賴 AI 完成所有學習任務。

在倫理與使用態度方面，學生普遍認同應學習如何正確使用 AI，並對 AI 可能產生錯誤資訊抱持一定程度的警覺，也認為直接繳交 AI 生成內容並不恰當，顯示學生對 AI 使用仍具備一定的學術倫理道德。

整體而言，本分析結果可以看出，生成式 AI 已逐漸成為部分學生學習過程中的輔助工具。學生大多肯定 AI 在學習上的便利性與幫助，但同時也對其使用方式與資訊正確性保持一定程度的注意與判斷。
\subsection{教師的觀點}\label{sec:教師觀點}

\subsubsection{填答教師基本背景}

本研究共調查十六位中小學教師，以探討教學現場對生成式人工智慧的看法。受訪者在任教階段以國中教師佔絕大多數，共計十三位，約佔百分之八十一點三；其餘三位為國小教師，其中包含一位特殊教育教師。學科分布呈現多元特徵，涵蓋科技與國文科各五位、英文或雙語學科三位，社會、自然及特教領域則各有一位。此分布顯示科技與語文類科目的教師對此議題展現出較高的參與度。就教學年資而言，受訪群體具備豐富的教學現場經驗，其中教學年資達十五年以上者共八位，十一至十五年者兩位，而初任教師（一至五年）則有六位。

\subsubsection{核心量化感受分析}

針對教師對生成式人工智慧的核心態度，本研究採用一至五分的李克特尺度進行量化評估。調查結果顯示，教師普遍肯定該技術在行政與備課上的生產力，在「提升教學效率」題項中，平均得分接近\textbf{四點五分}，多數教師給予四至五分的高分。然而，對於「有助於學生自主學習」之陳述，教師態度相對保留，平均得分約為三點六分，且評價呈現一分與三分的兩極化分布，顯示教師認為其成效高度取決於學生的自主學習能力。教學現場對於「降低學生思考能力」展現出集體焦慮，該題平均得分達\textbf{四點一分}，且有高達十二位教師給予四至五分。在「讓作業評量變得困難」方面，教師意見分歧，平均得分約為三點三分，評分從一分至五分皆有分布，這反映出評量難度與教師是否調整評量策略密切相關。最後，教師對於「需要具備人工智慧素養」與「人工智慧將改變未來教育模式」兩項陳述呈現高度共識，平均得分均達\textbf{四點七分以上}，絕大多數教師皆給予五分滿分。

\subsubsection{生成式人工智慧對教學現場之影響}

在開放式回饋中，教師指出生成式人工智慧的引入帶來了教師端正面賦能與學生端負面衝擊的強烈對比。主要改變可歸納為\textbf{三個核心面向}。首先是\textbf{教師備課與行政效率的提升}。多位科技、國小及語文科教師指出，該技術大幅減少了備課與撰寫行政文件的時間，使教師能將更多精力專注於教學設計與學生輔導。其次是\textbf{學生學習型態的兩極化轉變}。一方面，部分學生出現被動依賴與思考退化的現象，傾向直接尋求現成答案而減少了主動探究的過程，導致教師需花費更多心力辨識作業的原創性；另一方面，該技術也降低了創作與探究的門檻，使原本缺乏動力的學生更願意主動向系統提問，並能快速進行音樂或影片創作。最後是\textbf{教學互動模式的轉變}。由於資訊生成速度極快，教師角色正從單向的知識傳授者，轉變為引導者與檢驗者，必須具備更強的資訊辨識能力，以應對學生透過科技獲取的龐雜資訊。

\subsubsection{教學現場之因應策略與反思}

面對科技帶來的機遇與挑戰，教學現場的教師展現出務實的因應策略。在學習成效方面，教師認為該技術是強大的全天候學習夥伴，能即時解答疑惑，對低成就或性格內向的學生尤具輔助效果；但其帶來的資訊真偽判讀挑戰與速成依賴，亦嚴重影響學生批判是非的能力。為此，教師並未盲目禁止，而是採取\textbf{實體防禦與融入引導雙軌策略}。防守策略上，教師將重要寫作與評量限制在課堂內以紙筆或口說形式執行，以驗證學生的真實能力。引導策略上，語文與社會科教師刻意將其作為教學鷹架，例如要求學生分析系統修改句子的合理性，或引導學生批判系統整理的數據，藉此培養學生的批判性思考。關於課綱學習目標，教師普遍認為大原則不變，但教學重點應從死記背誦轉向理解、脈絡與判斷，閱讀理解、問題解決與資訊判讀素養等基礎能力在科技時代顯得更為關鍵。

\subsubsection{小結}

綜上所述，臺灣中小學教育現場正處於人工智慧衝擊的過渡期。教師在享受備課紅利的同時，也高度警覺學生因便利性而放棄思考的危機。親師生需建立正確觀念，將其定位為輔助工具，著重培養學生的素養與倫理。師資端自身亦需主動擁抱科技並提升專業知能，以掌握教學的平衡點。未來的教育政策與學校支持系統應著重於幫助教師發展引導學生批判科技內容的教案，並設計更具防弊功能且能引導高層次思考的多元評量方式，而非僅流於工具操作層面的研習，如此方能引導學生在科技時代奠定獨立思考與判斷是非的關鍵能力。

\section{文獻回顧：實徵研究的發現}\label{sec:文獻回顧}

為了解 AI 融入教育的影響，此部分針對10篇近期AI融入教育之實徵研究進行進一步分析與比較。考量第二部分調查結果，中小學師生最常使用之 AI 軟體為 GPT 與 Gemini，因此本部分以大型語言模型應用為核心，並納入與其功能相近的生成式 AI 或 AI 輔助程式學習工具（如 Claude、GitHub Copilot、iFLYTEK Spark、Google Colab 與 Gemini）之相關研究進行探討。此外，進一步將文獻回顧範疇限定於\textbf{程式學習領域}，以了解 AI 是否能協助學生解決電腦相關問題，研究對象涵蓋中小學生至大學生。

分析面向則包含實驗對象（人數與年齡）、AI型態、AI角色（例如如何融入程式學習）、探討變項，以及研究結果與討論等層面進行整理與比較。

\subsection{文獻名稱與編碼}\label{sec:文獻名稱}

\begingroup
\small
\setlength{\tabcolsep}{3pt}
\begin{longtable}{@{}p{0.16\textwidth} p{0.76\textwidth}@{}}
\toprule
編碼 & 文章名稱 \\
\midrule
\endfirsthead
\toprule
編碼 & 文章名稱 \\
\midrule
\endhead
\midrule
\multicolumn{2}{r}{續下頁} \\
\endfoot
\bottomrule
\endlastfoot
\cite{cao2023chatbots} & Cao, C. C., Ding, Z., Lin, J., \& Hoffgartner, F. (2023). AI Chatbots As multi-role pedagogical agents: Transforming engagement in CS education. \textit{arXiv Preprint} arXiv:2308.03992. \\
\cite{fan2025pair} & Fan, G., Liu, D., Zhang, R., \& Ran, L. (2025). The impact of AI-assisted pair programming on student motivation, programming anxiety, collaborative learning, and programming performance: a comparative study with traditional pair programming and individual approaches. \textit{International Journal of STEM Education}, 12(1), 16. \\
\cite{groothuijsen2024chatbots} & Groothuijsen, S., van den Beemt, A., Remmers, J. C., \& van Meeuwen, L. W. (2024). AI chatbots in programming education: Students' use in a scientific computing course and consequences for learning. \textit{Computers and Education: Artificial Intelligence}, 7, 100290. \\
\cite{llerena2024introductory} & Llerena-Izquierdo, J., Mendez-Reyes, J., Ayala-Carabajo, R., \& Andrade-Martinez, C. (2024). Innovations in introductory programming education: The role of AI with Google Colab and Gemini. \textit{Education Sciences}, 14(12), 1330. \\
\cite{omeh2025pbl} & Omeh, C. B., Ayanwale, M. A., Mnguni, L. E., \& Olelewe, C. J. (2025). Fostering programming skill and critical thinking through AI-assisted PBL integration. \textit{Journal of New Approaches in Educational Research,} 14(1), 22. \\
\cite{yan2025llm} & Yan, Y. M., Chen, C. Q., Hu, Y. B., \& Ye, X. D. (2025). LLM-based collaborative programming: impact on students' computational thinking and self-efficacy. \textit{Humanities and Social Sciences Communications}, 12(1), 1-12. \\
\cite{kohen2025integrating} & Kohen-Vacs, D., Usher, M., \& Jansen, M. (2025). Integrating generative AI into programming education: Student perceptions and the challenge of correcting AI errors. \textit{International Journal of Artificial Intelligence in Education}, 1-19. \\
\cite{yang2025chatgpt} & Yang, T. C., Hsu, Y. C., \& Wu, J. Y. (2025). The effectiveness of ChatGPT in assisting high school students in programming learning: evidence from a quasi-experimental research. \textit{Interactive Learning Environments}, 1-18. \\
\cite{tang2025enhancing} & Tang, B., Liang, J., Hu, W., \& Luo, H. (2025). Enhancing programming performance, learning interest, and self-efficacy: The role of large language models in middle school education. \textit{Systems}, 13(7), 555. \\
\cite{cubillos2025generative} & Cubillos, C., Mellado, R., Cabrera-Paniagua, D., \& Urra, E. (2025). Generative artificial intelligence in computer programming: Does it enhance learning, motivation, and the learning environment? \textit{IEEE Access.} \\
\end{longtable}
\endgroup

\subsection{實驗對象（人數、年齡）、AI型態與角色}\label{sec:實驗對象}

\begingroup
\small
\setlength{\tabcolsep}{3pt}
\begin{longtable}{@{}p{0.10\textwidth} p{0.09\textwidth} p{0.18\textwidth} p{0.15\textwidth} p{0.31\textwidth}@{}}
\toprule
文章 & 人數 & 年齡 & AI型態 & AI角色與功能 \\
\midrule
\endfirsthead
\toprule
文章 & 人數 & 年齡 & AI型態 & AI角色與功能 \\
\midrule
\endhead
\midrule
\multicolumn{5}{r}{續下頁} \\
\endfoot
\bottomrule
\endlastfoot
\cite{cao2023chatbots} & 200位 & 電腦科學碩士課程碩士生 & GPT-4 & 指導機器人、同儕機器人、職涯諮詢、情感支持 \\
\cite{fan2025pair} & 234位 & 大學生 & GPT-3.5 / Claude & 即時回饋、個人化協助、雙人程式設計即時建議、解釋概念 \\
\cite{groothuijsen2024chatbots} & 29位 & 機械工程與物理碩士生 & ChatGPT-3 / 3.5 & 解釋程式碼、解決方案、生成程式碼、檢查錯誤與除錯、概念理解、優化程式碼、解決數學問題 \\
\cite{llerena2024introductory} & 250位 & 程式設計課程大學生 & Google Colab & 個別化支持、演算法資源 \\
\cite{omeh2025pbl} & 62位 & 二年級電腦科學教育大學生 & Google Gemini & 即時回饋、生成演算法、優化程式碼、除錯任務、後設認知反思、錯誤解釋並建議、強化自主學習 \\
\cite{yan2025llm} & 85位 & 六和七年級的學生 & GPT-4 & 結構化鷹架、問題化鷹架、知識型鷹架 \\
\cite{kohen2025integrating} & 211名 & 大學生 & ChatGPT / GitHub Copilot & 學習夥伴、即時支援工具、學習鷹架、自動化助手 \\
\cite{yang2025chatgpt} & 153位 & 高二年級學生 & ChatGPT & 語法澄清、除錯、提供概念解釋 \\
\cite{tang2025enhancing} & 103位 & 七年級中學生 & iFLYTEK Spark & 個人化指導、即時回饋、智慧學習夥伴 \\
\cite{cubillos2025generative} & 40位 & 大學生 & Google Gemini 1.5 & 回答概念性和技術性問題 \\
\end{longtable}
\endgroup

\subsection{探討變因、研究結果與討論}\label{sec:探討變因}

\begingroup
\small
\setlength{\tabcolsep}{3pt}
\begin{longtable}{@{}p{0.09\textwidth} p{0.20\textwidth} p{0.29\textwidth} p{0.27\textwidth}@{}}
\toprule
代碼 & 探討變因 & 結果 & 討論 \\
\midrule
\endfirsthead
\toprule
代碼 & 探討變因 & 結果 & 討論 \\
\midrule
\endhead
\midrule
\multicolumn{4}{r}{續下頁} \\
\endfoot
\bottomrule
\endlastfoot
\cite{cao2023chatbots} & 探討AI回答的品質：準確性與流暢性、同理心與參與度、相關性與連貫性 & AI能夠提供準確且類似人類的回答；AI對於高認知水準的問題展現出更高的參與度與同理心；AI生成的回答與學生的問題相關且連貫 & AI 聊天機器人不僅能處理基礎問題，還能在學生進行更深入的探究式學習時，提供更有品質的互動和情感支持，進而達到更高的參與效果。 \\
\cite{fan2025pair} & 探討學生使用AI的：內在動機、程式設計焦慮、程式設計成效、協作學習認知、社會臨場感 & 內在動機：AI組顯著高於個人組、與人類組的成果相當；程式設計焦慮：AI組的焦慮減少幅度顯著大於個人組；程式設計成效：AI組的表現略優於個人組和人類組；社會臨場感：人類組的協作認知最高且顯著高於 AI 組 & AI 輔助配對程式設計顯著提高了學生的內在動機，並大幅降低了程式設計焦慮，效果與人類與人類配對相當，可能歸因於AI非評判性和即時回饋特性。儘管 AI 輔助編程改善了協作認知，但它無法達到人類與人類配對達到的協作深度和社會臨場感，強調了 AI 支援與人際互動的互補優勢。 \\
\cite{groothuijsen2024chatbots} & 學習成果、對科技的態度、教學活動的影響 & \textbf{學生和教師在 AI 影響學習方面的部分對立看法}：學生認為使用 ChatGPT 提升了他們的學習效果，讓他們寫出更簡潔的程式碼，並理解錯誤發生的原因。教師認為學生的程式作業整體品質不如往年，擔心學生過度依賴 ChatGPT，導致對學習的時間投入減少，而影響產出品質，僅對語法和語義學到了足夠的程度，對應用的學習不足。\textbf{學生和教師都同意使用 ChatGPT 對課程中的結對編程練習產生了負面影響}：學生只充當「驅動者」撰寫程式碼，並將 ChatGPT 用作「領航員」提供協助與反思，從而限制了他們與同儕之間的溝通和協作，教師觀察到團隊間的協作能力下降。 & 雖然 AI 機器人尚不能完全解決複雜的程式作業，但教學應更細緻地闡明哪些程式設計能力是學生必須掌握的，例如：更詳細的語法、語義、語用，以及哪些部分可以外包給 AI 工具。鑑於 AI 聊天機器人對協作的影響，必須進一步考慮是否以及如何維持結對編程原則。 \\
\cite{llerena2024introductory} & 程式設計概念理解、學習複雜主題的難易度、對主題的興趣、學習程式設計的動機、滿足期望、超越期望 & AI能幫助程式設計概念理解 87\%（55\% 非常同意，32\% 同意）；AI能幫助學習複雜主題 89\%（47\% 非常同意，42\% 同意）；AI能幫助提升對主題的興趣 91\%（51\% 非常同意，40\% 同意）；AI能提升學習程式設計的動機 86\%（46\% 非常同意，40\% 同意）；AI能提升滿足學生的期望 90\%（46\% 非常同意，44\% 同意）；AI能超越學生的期望 91\%（49\% 非常同意，42\% 同意） & 對基本概念如向量和陣列，以及控制、選擇和重複結構的設計的理解產生顯著影響，有助於認知發展的過渡，對基本概念理解更好的學生，也傾向於更容易理解複雜主題。 \\
\cite{omeh2025pbl} & 電腦程式設計能力、批判性思維能力、問題解決能力 & AI對程式設計能力有強烈的正向影響；AI 對批判性思維能力沒有顯著影響；AI 對問題解決能力沒有顯著影響 & AI能提升程式設計能力，但對批判性思維和問題解決能力未提升，學生對 AI 的過度依賴可能減少所需的認知努力。教學策略與學業能力之間存在顯著的交互作用，高學業能力的學生從程式設計教學中獲益更多。 \\
\cite{yan2025llm} & 運算思維、程式設計自我效能、認知負荷、心智負荷、心智努力 & AI組顯著提升運算思維；程式設計自我效能無顯著差異；AI組顯著降低認知負荷；AI組顯著降低心智負荷；心智努力無顯著差異 & AI能提升運算思維以更快地解決問題；AI組認知負荷顯著降低；教師不應只專注於讓學生獲得單一解決方案，而應著重於培養學生向AI提問並獲取答案的能力。 \\
\cite{kohen2025integrating} & 科技接受模型、任務類型 & AI組表現顯著低於傳統任務；使用AI時認知需求增加；AI知覺有用性高；AI能協助創意激發；強烈使用AI意願；對使用AI擔憂程度中等 & 雖然學生對 GenAI 抱持高度正面且熱情的看法，認為它能提升效率與創意；但在實際操作中，他們在面對 AI 生成的內容卻表現出明顯的困難。教育者在整合 GenAI 時，必須加強培養學生批判性評估與修正 AI 輸出內容的能力。 \\
\cite{yang2025chatgpt} & 心流體驗、自我效能、學習成就 & AI組全部顯著低於控制組 & 一些學生表示他們寧願選擇傳統方法，認為老師解釋事情比 ChatGPT 更清楚，並且可以迅速回應學生。學生描述模糊或過於專業時，ChatGPT 可能難以有效回應，導致學生花費大量時間澄清問題而非解決問題。AI可能無法迅速評估學習者的理解程度並調整挑戰水準，這可能會破壞學習平衡。由於AI能夠持續產出解決方案，學習者可能傾向於將學習成功歸因於 AI，從而削弱了他們從解釋自身表現中獲得的自我效能。 \\
\cite{tang2025enhancing} & 程式設計表現、學習興趣、自我效能 & 實踐分數：AI組實踐分數顯著高於傳統學習；認知興趣：AI組顯著高於傳統學習；自我效能：AI組顯著高於傳統學習 & 根據自我決定理論，AI透過開放式提問支持了學生的自主性，透過即時回饋增強了他們的能力感，並透過持續互動培養了關聯性，從而提升內在動機。互動鼓勵學生考慮多種解決方案，並探索超出規定課程的高級概念，刺激了程式設計知識的轉移和更深層次的理解。AI類似於一個隨時可用的動態鷹架，支援學生建立精通經驗和內在學習動機。 \\
\cite{cubillos2025generative} & 學習成效、感知能力、感知選擇/自主性、感知努力、感知壓力/緊張、感知實用性/價值 & 學習成效：無顯著差異；感知能力：無顯著差異；感知選擇/自主性：AI組顯著增加；感知努力：AI組顯著減少；感知壓力/緊張：AI組顯著減少；感知實用性/價值：無顯著差異 & AI組顯著提高了學生對感知自主性的認知，影片學習法顯著提高了自我效能，AI顯著減少了學生感知到的努力和壓力。AI組和影片這兩種方法各自具有獨特的優勢可以互補，教師可以考慮結合使用 AI組和教育影片。 \\
\end{longtable}
\endgroup

\subsection{(一) AI的角色}\label{sec:AI角色}

\begingroup
\small
\setlength{\tabcolsep}{3pt}
\begin{longtable}{@{}p{0.10\textwidth} p{0.18\textwidth} p{0.18\textwidth} p{0.18\textwidth} p{0.18\textwidth}@{}}
\toprule
\multicolumn{5}{l}{\textbf{AI角色類型}} \\
\midrule
項目 & 1. 導師與支持 & 2. 回饋與協助 & 3. 程式與除錯 & 4. 學習與合作 \\
\midrule
\endfirsthead
\toprule
\multicolumn{5}{l}{\textbf{AI角色類型}} \\
\midrule
項目 & 1. 導師與支持 & 2. 回饋與協助 & 3. 程式與除錯 & 4. 學習與合作 \\
\midrule
\endhead
\midrule
\multicolumn{5}{r}{續下頁} \\
\endfoot
\bottomrule
\endlastfoot
功能 & 概念澄清、職涯諮詢、鷹架（結構、問題）、情感支持 & 即時回饋、個別化內容/難度、加強概念理解、解釋程式碼、語法澄清 & 解決/生成程式碼、支援演算法、優化程式碼、除錯任務、解決複雜問題 & 結對程式設計、協助同儕配對、促進後設認知反思、使用者介面增強提示、強化自主學習、參與度監測 \\
文章 & \cite{cao2023chatbots,yan2025llm,kohen2025integrating,tang2025enhancing,cubillos2025generative} (5篇) & \cite{fan2025pair,llerena2024introductory,omeh2025pbl,kohen2025integrating,yang2025chatgpt,tang2025enhancing} (6篇) & \cite{fan2025pair,groothuijsen2024chatbots,llerena2024introductory,omeh2025pbl,kohen2025integrating,yang2025chatgpt} (6篇) & \cite{fan2025pair,omeh2025pbl,kohen2025integrating} (3篇) \\
\end{longtable}
\endgroup

綜整十篇文獻後發現，AI在教育情境中的角色大致可歸納為四種類型。第一類為導師與支持者，主要協助學生進行概念澄清、提供職涯諮詢、回應結構性問題以及情感支持；第二類為回饋與協助者，著重於提供即時回饋、依學生需求調整內容與難度、加強概念理解，以及協助程式碼與語法修正；第三類為程式與除錯助手，能協助生成與修正程式碼、支援演算法設計、優化程式效能及解決複雜技術問題；第四類則為學習與合作夥伴，透過協助程式設計、促進同儕配對、引導後設認知反思、提供適時提示以及監測學習參與度，支持學生的自主學習與合作學習歷程。整體而言，文獻顯示AI的角色已不僅限於知識提供工具，而是逐漸朝向兼具教學支持、學習回饋、技術協助及學習夥伴等多元功能發展，其中以回饋協助與程式除錯相關功能最為常見，反映AI在提升學習效率與解決學習問題方面的重要價值。

\subsection{(二) 優缺點比較}\label{sec:優缺點比較}

\begingroup
\small
\setlength{\tabcolsep}{2pt}
\begin{longtable}{@{}p{0.10\textwidth} p{0.115\textwidth} p{0.115\textwidth} p{0.115\textwidth} p{0.115\textwidth} p{0.115\textwidth} p{0.115\textwidth}@{}}
\toprule
\multicolumn{7}{l}{\textbf{優勢}} \\
\midrule
類別 & 成效提升 & 焦慮下降 & 幫助理解 & 降低認知負荷 & 提升認知參與度 & 提升自我效能 \\
\midrule
\endfirsthead
\toprule
\multicolumn{7}{l}{\textbf{優勢}} \\
\midrule
類別 & 成效提升 & 焦慮下降 & 幫助理解 & 降低認知負荷 & 提升認知參與度 & 提升自我效能 \\
\midrule
\endhead
\midrule
\multicolumn{7}{r}{續下頁} \\
\endfoot
\bottomrule
\endlastfoot
文章 & \cite{fan2025pair,groothuijsen2024chatbots,llerena2024introductory,omeh2025pbl,kohen2025integrating,tang2025enhancing} & \cite{fan2025pair} & \cite{cao2023chatbots,groothuijsen2024chatbots,llerena2024introductory,tang2025enhancing} & \cite{yan2025llm} & \cite{cao2023chatbots,kohen2025integrating} & \cite{yan2025llm,tang2025enhancing,cubillos2025generative} \\
\end{longtable}
\endgroup

\begingroup
\small
\setlength{\tabcolsep}{2pt}
\begin{longtable}{@{}p{0.10\textwidth} p{0.115\textwidth} p{0.115\textwidth} p{0.115\textwidth} p{0.115\textwidth} p{0.115\textwidth} p{0.115\textwidth}@{}}
\toprule
\multicolumn{7}{l}{\textbf{劣勢}} \\
\midrule
類別 & 協作學習認知低 & 低社會臨場感 & 協作能力下降 & 降低學習心流 & 任務表現下降 & 自我效能下降 \\
\midrule
\endfirsthead
\toprule
\multicolumn{7}{l}{\textbf{劣勢}} \\
\midrule
類別 & 協作學習認知低 & 低社會臨場感 & 協作能力下降 & 降低學習心流 & 任務表現下降 & 自我效能下降 \\
\midrule
\endhead
\midrule
\multicolumn{7}{r}{續下頁} \\
\endfoot
\bottomrule
\endlastfoot
文章 & \cite{fan2025pair,groothuijsen2024chatbots} & \cite{fan2025pair} & \cite{groothuijsen2024chatbots} & \cite{yang2025chatgpt} & \cite{kohen2025integrating} & \cite{yang2025chatgpt} \\
\end{longtable}
\endgroup

綜整十篇文獻後發現，AI融入教學整體上展現出多項正向效益。在優勢方面，最常被提及的是\textbf{提升程式設計成效}，相關研究指出AI能協助學習者理解程式概念、改善程式品質及提升問題解決能力\cite{fan2025pair,groothuijsen2024chatbots,llerena2024introductory,omeh2025pbl,kohen2025integrating,tang2025enhancing}。此外，AI亦有助於\textbf{理解複雜問題}\cite{cao2023chatbots,groothuijsen2024chatbots,llerena2024introductory,tang2025enhancing}、\textbf{提升自我效能}\cite{yan2025llm,tang2025enhancing,cubillos2025generative}及\textbf{提高認知參與度}\cite{cao2023chatbots,kohen2025integrating}，使學習者在面對學習任務時更具信心與投入感。部分研究亦指出，AI能降低程式設計焦慮\cite{fan2025pair}及認知負荷\cite{yan2025llm}，進而提升學習效率與學習體驗。

然而，文獻亦揭示AI融入教學可能帶來許多負面影響。在協作學習情境中，部分研究發現學生可能因過度依賴AI而出現\textbf{協作學習認知降低}\cite{fan2025pair,groothuijsen2024chatbots}及\textbf{協作能力下降}\cite{groothuijsen2024chatbots}的情形；同時，AI的介入亦可能降低學習者的學習心流\cite{yang2025chatgpt}與社會臨場感\cite{fan2025pair}，影響人際互動與合作品質。此外，有研究指出AI目前仍難以即時且精準地評估學習者的理解程度\cite{yang2025chatgpt}，並可能造成部分學生自我效能下降\cite{yang2025chatgpt}。

\cite{kohen2025integrating,yang2025chatgpt}兩篇研究皆指出，雖然學生普遍肯定AI在學習上的便利性與支持功能，但其應用仍可能帶來若干負面影響。首先，\textbf{學生在修正大型語言模型（LLM）生成的程式碼時，往往面臨比傳統程式任務更大的挑戰}，原因在於AI產生的錯誤通常較為隱晦且複雜，需要較高層次的程式知識與除錯能力才能辨識與修正。此外，部分學生認為教師的解說較AI清楚且更具即時互動性，\textbf{當學生提出模糊或過於專業的問題時，AI可能無法準確理解需求，反而增加釐清問題所耗費的時間}。AI難以像教師一樣即時評估學習者的理解程度並調整教學內容與難度，可能出現內容過於簡單或過於複雜的情況，影響學習效果\cite{yang2025chatgpt}。

綜合而言，兩篇研究顯示AI雖能提供有效的學習支援，但若缺乏適當引導與使用策略，仍可能對學生的自主學習、問題解決能力及學習信心產生負面影響\cite{kohen2025integrating,yang2025chatgpt}。AI融入教學雖能有效提升學習成效與學習支持，但是否會讓學生過度依賴及對協作互動的影響仍值得進一步探討，未來應在發揮AI輔助功能與維持學生自主學習及人際互動之間取得平衡。

\section{組員立場：我們的觀點}\label{sec:組員立場}

\subsection{慶霖（贊成）}\label{sec:慶霖}

我認為AI應融入教學中。AI的發展能分擔教學現場許多重複性的工作，並為個別化教學帶來進一步的突破。雖然目前AI仍存在不少問題，且尚未建立一套具體明確的使用指引，但若因此拒絕使用，便無法深入了解其潛力與局限。任何教育科技皆有其優缺點。隨著時間推移，人們能透過實踐經驗與實徵研究，更清楚如何將不同科技應用於合適的學科。目前生成式AI融入教學仍處於探索階段，需要更多研究與教學設計來實證其成效。AI有望帶來比以往更佳的教學成效，但我們也必須投入心力防範學生過度依賴。從教學現場來看，可先從減輕教師的重複性工作著手。面對臺灣教師工作量繁重的現況，AI能協助教師減少備課與批改作業的負擔。此外，在教學設計上應高度重視實徵研究指出的問題，避免出現「AI幫學生完成了任務，學生自身卻毫無所得」的困境。

\subsection{吟貞（反對）}\label{sec:吟貞}

我認為AI不應融入教學中。雖然本次問卷調查結果顯示，多數學生自認不會過度依賴AI，但在實際觀察中，我發現身邊許多同學已然出現過度依賴的現象，遇到問題時直覺地求助AI，而不再嘗試自行思考與解題。長此以往，缺乏主動思考的習慣將導致思維能力逐步下降，且思維能力一旦退化便難以恢復。當學生逐漸喪失獨立判斷與批判性分析的能力，教育原本應當培養的核心素養也將淪為空談。因此，在尚未建立有效的規範與防護機制之前，貿然將AI融入教學，反而可能削弱學生最應具備的基礎能力。

\subsection{詠傑（贊成）}\label{sec:詠傑}

我認為AI應融入教學中。在現行教育體制下，單一教師往往需要同時面對數十位學生，難以逐一回應每位學生的個別問題；若能善用AI，學生便能更快速地獲得即時解答，降低等待的時間成本。對於具備自學意願的學生而言，AI能大幅提升其學習效率與自主探究能力，讓自主學習不再受限於教師資源的分配。此外，AI能協助教師更高效地備課與規劃課程，將原本耗費在重複性庶務上的時間釋放出來，轉而投入於更具價值的教學設計與個別輔導。最後，隨著AI在各行各業的普及，積極接觸與使用AI能讓我們更深入了解如何高效運用這項工具。這不僅符合未來趨勢，更是面對AI時代不可或缺的素養。

\subsection{駿甫（贊成）}\label{sec:駿甫}

我認為AI應融入教學中。將AI納入包含教育在內的各個領域已是不可逆的趨勢，與其被動抗拒，不如主動思考如何善用。此外，現階段並無有效方法能完全防範學生使用AI。既然無法禁止，就應正視此一現實，將AI納入教學框架之中。透過明確的規範與引導，學生能在安全的環境下學習如何正確使用AI，而非在缺乏指引的情況下自行摸索，進而造成更不可控的風險。

\section{結論}\label{sec:結論}

綜合教學現場調查與文獻回顧，我們發現：

\subsection{(一) AI帶來哪些教育現場的效益？}\label{sec:結論一}

\begin{itemize}
    \item \textbf{學生端}：提供更具個別化的學習環境，如個別化回饋、解決能力差異的配對問題、個別化的指導等。正因為學生能夠對AI提出任何提問，而AI也能夠給出回覆，不僅能夠對學生給予個別化的回饋，讓學生能夠在任何時間點提出任何問題，且重複提問直到真正理解為止，也能夠讓AI換一種方式來說明，以配合學生的理解程度差異，這樣的學習支持有效地改善了過去教育科技的侷限，過去的教育科技多侷限於某個學科或單元，且使用方式較為侷限，學生可操作的空間較小，而生成式AI相較於其他教育科技具有更大的彈性，讓學生可以根據自己所需來使用AI工具，為個別化學習提供更大的可能性。

    \item \textbf{教師端}：AI幫助教師在備課或作業批改上有更多協助，在教學現場端有許多重複性的工作，如作業批改、成績計算等，若AI能夠幫助教師完成這些部分，教師則能夠更專注於教學設計或提供學生更個別化的學習建議，以減輕教師的負擔。
\end{itemize}

此外，相關研究指出AI能夠帶來提升學習成效、降低焦慮、幫助概念理解、降低認知負荷、提升認知參與度、提升自我效能等幫助，顯示了其在各個方面為學習帶來的助益。

\subsection{(二) AI為教育帶來哪些擔憂？}\label{sec:結論二}

從現場調查與研究回顧可見，許多教師擔憂學生使用AI產生的依賴問題，即便學生透過AI可以完成更多的學習任務，並減輕學習壓力，然而許多研究指出，在AI融入教育的情況下，更應注重學生\textbf{批判性思考能力}，因學生要能夠判斷AI提供的內容是否正確，就必須具備批判AI提供的內容的能力，而非全然相信，尤其針對此種非針對教育而設計的大型語言模型，產出的內容往往可能有誤，若學生不具備批判性思考的能力，則最終可能學習到錯誤的知識。此外，在使用方法上，若學生僅向AI索取答案，跳過了思考的部分，則會讓過去的評量標準失去其功能性，即便學生繳交出正確答案，仍然無法確保其真正理解，因此在AI融入教學的設計中，其評量方式需要做出改變。

\subsection{(三) AI是否影響教育的目標？}\label{sec:結論三}

十二年國教課程綱要總綱中指出：

\begin{quote}
「核心素養」是指一個人為適應現在生活及面對未來挑戰，所應具備的知識、能力與態度。

「國家課程規範歷經數次中小學課程標準修訂，務求課程修訂能與時俱進。以培養健全國民為宗旨，為我國人才培育奠定良好基礎。」
\end{quote}——\cite{moe2021curriculum}

課綱的修訂會隨著時代變遷而有所修訂，目前AI已經取代部分工作時，顯示了AI已開始對社會造成改變，人們已經不可避免地會接觸並使用AI，因此教育部在未來的課綱修訂中，應考量如何使學生具備使用AI的能力，並且同時兼具在「自主行動」、「溝通互動」、「社會參與」等面向的學習目標，AI融入皆可能會對這些面向產生衝擊，例如：在「自主行動」面向中

\begin{quote}
「應培養學生具備問題理解、思辨分析、推理批判的系統思考與後設思考素養，並能行動與反思，有效處理及解決生活、生命問題。」
\end{quote}——\cite{moe2021curriculum}

若使用AI時直接獲取答案，則無法培養問題理解、思辨分析、推理批判等能力，因此未來在AI融入教學上應謹慎評估學生是否具備這些能力，且應培養學生對AI的了解，以及使用AI時具備批判性思考的能力。

\subsection{(四) AI會做這些事情，我們還要學習嗎？}\label{sec:結論四}

\begin{quote}
「教育是為適應現在生活及面對未來挑戰，成為具有社會適應力與應變力的終身學習者。」
\end{quote}——\cite{moe2021curriculum}

AI幫助我們的生活更便利，讓我們可以不用付出相等的努力就完成過去無法完成的事情，然而這不代表著我們就能夠因此不去學習，而是應該重新思考在AI時代下，我們應該學習的事情是什麼。隨著AI快速發展與變遷，有必要不斷地檢核教育的目標與方法，以適應AI對社會帶來的改變。人的時間及力氣是有限的，應該進一步思考在未來的課綱中哪些目標應該調整或加入，高階能力的培養起始於基礎能力，因此不應因為AI能夠取代這些基礎能力，就認為培養這些基礎能力並不重要，可以重新調整這些課程設計上的比重，以培養學生在各項指標上的能力。

\section{未來建議}\label{sec:未來建議}

\subsection{(一) 課綱的調整}\label{sec:未來建議一}

面對AI時代，十二年國教課綱應重新思考並納入AI作為一門學科或是AI融入教學的設計，並將其納入電腦課程或是強化其他課程在批判性思考的培養，以健全學生面對AI變遷的能力，並考量到AI普及化下教育目標應如何調整，確保學生有足夠的能力來應對社會發展。

\subsection{(二) AI融入教學設計}\label{sec:未來建議二}

在AI融入教案的設計上，應重新去檢視AI如何融入教學，以及過去評量方法是否還能夠有效評估學生的能力，在教師端也應特別去設計能夠檢核學生是否真正學會的環節，可透過讓學生記錄AI使用的方法、批判性思考環節等，避免學生過度依賴AI的問題。

\subsection{(三) 國家的AI使用指引}\label{sec:未來建議三}

面對生成式 AI 對教育體系的衝擊，臺灣亟需整合國際政策經驗，制定具體且具可操作性的國家級使用指引。本報告結合跨國政策比較與本土調查，提出以下五大核心政策建議：

第一，\textbf{建立分齡分層的 AI 素養教育體系}。我國應以教育部現行的《我和 AI 一起學》與《駕馭 AI，洞察未來》手冊為基礎\cite{govtw2024aiguide}，結合聯合國教科文組織（UNESCO）的「理解、應用、創造」三階段漸進模式\cite{unesco2024competency}，建構系統化素養指標。國小階段著重於辨識 AI 幻覺與隱私保護；國高中階段強調批判性查核與共學；高等教育則專注於高階提示詞工程與倫理責任。

第二，\textbf{重構評量機制，轉向過程驗證}。為避免如美國普渡大學因 AI 偵測軟體誤判而引發的校園信任危機\cite{arun2026ai}，評量機制必須從「結果導向」轉為「過程導向」。各級學校應減少回家作業的分數比重，並效法全國中小學科展的口頭答辯機制\cite{nehs2025science}，透過即席問答與實作演示，確認學生真正內化知識，而非僅是複製 AI 的產出。

第三，\textbf{落實強制揭露與使用者追責制度}。在學術與競賽中，應確立「創意歸屬人類，勞務外包 AI」的原則。參考新加坡國立大學的當責要求\cite{nus2024integrity}與旺宏科學獎的作法\cite{mxic2025competition}，強制要求學生在提交作業時，必須完整保留並揭露與 AI 協作的提示詞軌跡（Prompt Logs），且使用者須對所有產出內容負起最終責任。

第四，\textbf{推動教師賦能與系統性培訓}。借鑑英國教育部的雙軌策略\cite{dfe2023genai}與新加坡高達百分之七十五的教師 AI 使用率\cite{moe2025singapore}，國家應提供系統性的教師培訓，協助教師使用 AI 工具減輕行政與備課負擔，並提升教師設計 AI 融入課程與評估學生 AI 使用軌跡的專業能力。

第五，\textbf{強化風險評估與資料隱私保護}。在推動 AI 融入教學的同時，必須建立嚴格的安全防護網。參考英國教育部對網路監控與提示詞可見性的要求\cite{dfe2023genai}，以及日本文部科學省對學生思考力低下的警示\cite{mext2024genai}，學校在導入 AI 工具前應進行全面的風險評估，確保學生隱私安全，並引導學生維持主體思考，避免盲目接受 AI 答案。

\section{組員分工}\label{sec:組員分工}

\begin{longtable}{ll}
\toprule
成員 & 分工 \\
\midrule
陳慶霖 & 摘要、問卷設計與發放、文獻回顧與分析、結論、未來建議 \\
鍾詠傑 & 議題簡介、現況分析、未來建議 \\
林吟貞 & 學生態度分析 \\
林駿甫 & 教師觀點分析 \\
\bottomrule
\end{longtable}

\addcontentsline{toc}{section}{參考文獻}
\printbibliography[segment=\the\value{refsegment},title=參考文獻,heading=subbibliography]
